% ECONOMICS_PROBLEM: {"id": "F16-522", "title": "Trade shocks, inequality and local adjustment", "jel_code": "F16", "source_id": "OP-0083"}
% !TeX program = xelatex
\documentclass[11pt,a4paper]{article}
\usepackage[margin=23mm]{geometry}
\usepackage{fontspec}
\setmainfont{Latin Modern Roman}
\usepackage{amsmath,amssymb,mathtools}
\usepackage{xcolor,enumitem,booktabs,longtable,array,xurl,hyperref}
\definecolor{muted}{HTML}{53636C}
\hypersetup{colorlinks=true,urlcolor=blue,linkcolor=blue}
\setlength{\parindent}{0pt}
\setlength{\parskip}{5pt}
\newcommand{\R}{\mathbb R}
\newcommand{\E}{\mathbb E}
\newcommand{\N}{\mathbb N}
\newcommand{\ind}{\mathbf 1}
\newcommand{\Var}{\operatorname{Var}}
\newcommand{\Cov}{\operatorname{Cov}}
\newcommand{\supp}{\operatorname{supp}}
\DeclareMathOperator*{\argmin}{arg\,min}
\DeclareMathOperator*{\argmax}{arg\,max}
\newcommand{\entrymeta}[3]{{\sffamily\small\color{muted}\textbf{\texttt{#1}}\quad|\quad #2\par #3\par}}
\newcommand{\Problemref}[1]{\texttt{#1}}
\newcommand{\JELref}[2]{\texttt{#2} (JEL \texttt{#1})}
\begin{document}
\section*{Trade shocks, inequality and local adjustment}
\textbf{Problem ID:} \texttt{F16-522}\quad \textbf{JEL:} \texttt{F16}\par
% BEGIN ECONOMICS STATEMENT
\label{problem:OP-0083}
\label{atlas:prob:T3}

\noindent\textbf{Atlas ID:} T3; \textbf{original number:} 83; \textbf{field:} International trade. \textbf{Classification:} Editorial research formulation (theoretical/empirical); not a certified unsolved theorem.

\subsection*{Definitions and assumptions}
Fix regions $r=1,\ldots,R$, workers $h$ with initial region $r(h)$, and a finite horizon $H$. Let $q=(q_0,\ldots,q_H)$ describe an externally specified national trade-shock path. A dynamic spatial model specifies worker location and skill choices, housing supply, firm hiring and entry, migration costs, and transfer rules. Agents choose feasible actions subject to budgets and housing and labor markets clear. Include an equilibrium selection because congestion, housing and entry may yield multiplicity. Let $Y_{ht}(q)$ denote earnings or employment under the entire shock path, and require finite first moments.

\subsection*{Formal research question}
Identify the initial-residence impulse response
\[
 \Delta_r(t)=\mathbb E[Y_{ht}(q^1)-Y_{ht}(q^0)\mid r(h)=r],
 \qquad t=0,\ldots,H,
\]
and its decomposition into migration, within-region wages, employment and worker composition. For welfare, use a separate discounted utility comparison that includes housing costs, migration and transfers. Characterize observationally equivalent combinations of mobility costs, firm adjustment and skill accumulation. Determine which longitudinal worker, firm and housing measurements or policy interventions identify those mechanisms and predict recovery after shocks outside the estimation episode.

\subsection*{Interpretation and scope}
Using initial residence avoids mechanically interpreting migration-induced composition changes as incumbent-worker recovery. Potential outcomes depend on the full shock vector, allowing interference across regions through migration and trade. A shift-share exposure is a measurement device, not itself an exogenous treatment. An instrument based on foreign import growth requires a stated exclusion restriction and identification argument; global demand or technology shocks may violate it. If only those assumptions support a local effect, do not label the estimate a universal national welfare effect. Define persistence by a horizon and tolerance, or separately impose conditions allowing an infinite-horizon limit. Any extrapolation to other countries or shocks requires explicit invariance assumptions for preferences, technologies and mobility constraints.

\subsection*{Source audit and status}
[S1] documents historical US local effects, [S2] reviews adjustment and supplies atlas link [42], and [S3] studies Brazilian regional dynamics. These provide context and evidence in particular settings. The claim that every affected region remains scarred for decades is not warranted. The present question concerns mechanisms, transportability and prediction, not an unanswered existence theorem.

\subsection*{References}

\begin{enumerate}[label={[S\arabic*]},leftmargin=*]

\item David H. Autor, David Dorn, and Gordon H. Hanson (2013). \emph{The China Syndrome: Local Labor Market Effects of Import Competition in the United States}. American Economic Review 103(6), 2121–2168. \url{https://doi.org/10.1257/aer.103.6.2121}.
\textit{Locator and role:} Publisher or author abstract / bibliographic record; Local labor-market import-competition estimates in the historical US setting. Provides background or a benchmark result; does not establish that the editorial question remains unresolved in 2026.

\item David H. Autor, David Dorn, and Gordon H. Hanson (2016). \emph{The China Shock: Learning from Labor-Market Adjustment to Large Changes in Trade}. Annual Review of Economics 8, 205–240. \url{https://doi.org/10.1146/annurev-economics-080315-015041}.
\textit{Locator and role:} Publisher or author abstract / bibliographic record; Review of persistent labor-market adjustment; also resolves the atlas link [42]. Provides background or a benchmark result; does not establish that the editorial question remains unresolved in 2026.

\item Rafael Dix-Carneiro and Brian K. Kovak (2017). \emph{Trade Liberalization and Regional Dynamics}. American Economic Review 107(10), 2908–2946. \url{https://doi.org/10.1257/aer.20161214}.
\textit{Locator and role:} Publisher or author abstract / bibliographic record; Regional adjustment following Brazilian trade liberalization. Provides background or a benchmark result; does not establish that the editorial question remains unresolved in 2026.

\end{enumerate}

\subsection*{Common conventions: Source mathematical conventions}

Unless an entry states otherwise, agent and firm sets are finite. State and action spaces are measurable, strategies and outcome maps are measurable, probability laws are specified, and expectations exist for the targets under discussion. Infinite-horizon discount factors lie in $(0,1)$ and discounted objectives are finite. Norms, tolerances and complexity input sizes must be specified for computational comparisons. Constants or primitives that appear locally are inputs, not empirical facts asserted by this catalogue.

\textbf{Identification.} A statistical model $m\in\mathcal M$ implies an observable law $P_m$ and target $\theta(m)$. At law $P$, its identified set is
\[
\Theta(P;\mathcal M)=\{\theta(m):m\in\mathcal M,\ P_m=P\}.
\]
Point identification holds when this set is a singleton, or equivalently when $P_{m_1}=P_{m_2}$ implies $\theta(m_1)=\theta(m_2)$ for all compatible models. Sharp bounds mean bounds attained, or approached when the boundary is not attained, by compatible models. Writing this set is a definition, not a solution: the substantive task is to establish informative restrictions and derive or estimate the set. An unrestricted-confounding model may give only trivial bounds.

\textbf{Inference.} Identification, estimability and valid uncertainty quantification are separate questions. Fix $\alpha\in(0,1)$. A confidence procedure $C_n$ over a declared law class $\mathcal P$ has asymptotic uniform coverage at level $1-\alpha$ when
\[
\liminf_{n\to\infty}\inf_{P\in\mathcal P}\Pr_P\{\theta(P)\in C_n\}\geq1-\alpha.
\]
For set-identified targets the coverage event must be defined separately, for example coverage of the full identified set. If an entry uses identification-indexed models instead of uniquely indexed laws, the same coverage convention is applied over those models. Pointwise limits do not establish uniform validity, and asymptotic validity does not guarantee finite-sample precision. Sample sizes, dependence structures and weak-identification sequences are part of each application.

\textbf{Counterfactuals.} $Y_i(a)$ is a potential outcome under a specified intervention, including its timing, implementation and versions. It is not $Y_i$ conditional on voluntarily choosing $a$. A full intervention vector permits interference; shorthand scalar treatment notation does not silently impose no interference. Assignment, selection, attrition, anticipation and measurement must be specified. A claim about an instrument's local effect is not automatically a population policy effect.

\textbf{Economic equilibrium.} A counterfactual model includes best responses, constraints, feasibility, market clearing, state transitions and expectations consistent with the stated information. When several equilibria exist, either a declared selector is an input or the target is set-valued. Comparisons across policy regimes must use the stated selection convention. Reduced-form relationships are not presumed invariant after an equilibrium-changing intervention.

\textbf{Optimization and welfare.} Optimization is over the stated feasible set. If attainment is not established, $\sup$ or $\inf$ and $\varepsilon$-optimal choices replace an asserted maximizer or minimizer. For example, with value $V=\sup_{a\in\mathcal A}W(a)<\infty$, an $\varepsilon$-optimal set is $\{a\in\mathcal A:W(a)\geq V-\varepsilon\}$ for $\varepsilon>0$. Benefits, costs, risk and health or environmental outcomes can be aggregated only using declared units and conversion weights. Interpersonal utility weights, the treatment of fiscal transfers and foreign profits, discounting, and the behavioral welfare criterion are explicit inputs. A comparative-static derivative additionally requires existence and appropriate differentiability of the selected counterfactual.

\textbf{Sources.} Local labels [S1], [S2], and so forth restart in every question. Locators identify the relevant abstract, model, section or result at the level actually checked; a bibliographic or abstract-level verification is not represented as inspection of an entire proof. Working papers are distinguished from later publications and changing versions. Source summaries are paraphrased. All URL checks and editorial formulations refer to the audit date stated above.
% END ECONOMICS STATEMENT
\end{document}
